% Options for packages loaded elsewhere
\PassOptionsToPackage{unicode}{hyperref}
\PassOptionsToPackage{hyphens}{url}
\documentclass[
]{article}
\usepackage{xcolor}
\usepackage{amsmath,amssymb}
\setcounter{secnumdepth}{-\maxdimen} % remove section numbering
\usepackage{iftex}
\ifPDFTeX
  \usepackage[T1]{fontenc}
  \usepackage[utf8]{inputenc}
  \usepackage{textcomp} % provide euro and other symbols
\else % if luatex or xetex
  \usepackage{unicode-math} % this also loads fontspec
  \defaultfontfeatures{Scale=MatchLowercase}
  \defaultfontfeatures[\rmfamily]{Ligatures=TeX,Scale=1}
\fi
\usepackage{lmodern}
\ifPDFTeX\else
  % xetex/luatex font selection
\fi
% Use upquote if available, for straight quotes in verbatim environments
\IfFileExists{upquote.sty}{\usepackage{upquote}}{}
\IfFileExists{microtype.sty}{% use microtype if available
  \usepackage[]{microtype}
  \UseMicrotypeSet[protrusion]{basicmath} % disable protrusion for tt fonts
}{}
\makeatletter
\@ifundefined{KOMAClassName}{% if non-KOMA class
  \IfFileExists{parskip.sty}{%
    \usepackage{parskip}
  }{% else
    \setlength{\parindent}{0pt}
    \setlength{\parskip}{6pt plus 2pt minus 1pt}}
}{% if KOMA class
  \KOMAoptions{parskip=half}}
\makeatother
% definitions for citeproc citations
\NewDocumentCommand\citeproctext{}{}
\NewDocumentCommand\citeproc{mm}{%
  \begingroup\def\citeproctext{#2}\cite{#1}\endgroup}
\makeatletter
 % allow citations to break across lines
 \let\@cite@ofmt\@firstofone
 % avoid brackets around text for \cite:
 \def\@biblabel#1{}
 \def\@cite#1#2{{#1\if@tempswa , #2\fi}}
\makeatother
\newlength{\cslhangindent}
\setlength{\cslhangindent}{1.5em}
\newlength{\csllabelwidth}
\setlength{\csllabelwidth}{3em}
\newenvironment{CSLReferences}[2] % #1 hanging-indent, #2 entry-spacing
 {\begin{list}{}{%
  \setlength{\itemindent}{0pt}
  \setlength{\leftmargin}{0pt}
  \setlength{\parsep}{0pt}
  % turn on hanging indent if param 1 is 1
  \ifodd #1
   \setlength{\leftmargin}{\cslhangindent}
   \setlength{\itemindent}{-1\cslhangindent}
  \fi
  % set entry spacing
  \setlength{\itemsep}{#2\baselineskip}}}
 {\end{list}}
\usepackage{calc}
\newcommand{\CSLBlock}[1]{\hfill\break\parbox[t]{\linewidth}{\strut\ignorespaces#1\strut}}
\newcommand{\CSLLeftMargin}[1]{\parbox[t]{\csllabelwidth}{\strut#1\strut}}
\newcommand{\CSLRightInline}[1]{\parbox[t]{\linewidth - \csllabelwidth}{\strut#1\strut}}
\newcommand{\CSLIndent}[1]{\hspace{\cslhangindent}#1}
\setlength{\emergencystretch}{3em} % prevent overfull lines
\providecommand{\tightlist}{%
  \setlength{\itemsep}{0pt}\setlength{\parskip}{0pt}}
\usepackage{bookmark}
\IfFileExists{xurl.sty}{\usepackage{xurl}}{} % add URL line breaks if available
\urlstyle{same}
\hypersetup{
  pdftitle={Gators: Machine Learning Preprocessing and Feature Engineering on Polars with ONNX Export},
  hidelinks,
  pdfcreator={LaTeX via pandoc}}

\title{Gators: Machine Learning Preprocessing and Feature Engineering on
Polars with ONNX Export}
\author{}
\date{15 September 2026}

\begin{document}
\maketitle

\subsection{Summary}\label{summary}

Data preprocessing and feature engineering are executed far more often
than model training: every serving request re-runs the same
transformations that were fit once during training. \texttt{Gators} is
an open-source Python library that provides 107 transformers, organized
into 10 categories (data cleaning, outlier clipping, categorical
encoding -- including CatBoost-style target encoding (Prokhorenkova et
al. 2018) -- numeric/string/datetime feature generation, missing-value
imputation, discretization, scaling, and feature selection), implemented
on top of the \texttt{Polars} DataFrame engine (Vink and Polars
Contributors 2020) and exposed through a
\texttt{scikit-learn}-compatible \texttt{fit}/\texttt{transform} API
(Pedregosa et al. 2011). A fitted \texttt{Gators} pipeline can be
chained with \texttt{gators.pipeline.Pipeline} and exported, as a single
object, to a validated ONNX graph (ONNX Contributors 2019a) for
runtime-agnostic, language-agnostic inference --- without
hand-maintaining a second implementation of the same preprocessing logic
for production. The library is tested with 2,323 unit and
ONNX-conversion tests and ships a reproducible benchmark suite comparing
it against \texttt{scikit-learn} and \texttt{feature-engine} (Galli
2021) across row counts, feature counts, missingness, cardinality, and
thread configurations.

\subsection{Statement of need}\label{statement-of-need}

Preprocessing code is frequently treated as disposable glue rather than
a first-class part of a machine learning system, despite being both a
major source of compute cost on large tabular datasets and a major
source of production failures. Sculley et al. (Sculley et al. 2015)
identify data-pipeline glue code as one of the largest sources of
technical debt in ML systems, and Breck et al. (Breck et al. 2017) note
that \emph{training/serving skew} --- a mismatch between how a feature
is computed at training time versus serving time --- is one of the most
common causes of silent production degradation. The standard mitigations
are either (a) maintaining two implementations of the same feature logic
(a notebook version in \texttt{pandas} and a hand-written production
version in a second language), and testing the two for parity, or (b)
tying the preprocessing logic to a single serving runtime. Neither is
satisfying: (a) is a recurring maintenance and correctness burden, and
(b) trades away portability. \texttt{Gators} targets researchers and
practitioners who need (i) a wide, consistent catalogue of preprocessing
operations that composes with existing \texttt{scikit-learn} tooling,
(ii) an implementation fast enough that preprocessing is not the
bottleneck in a research iteration loop, and (iii) a way to take a
fitted pipeline to production inference without re-implementing it,
while being explicit about which parts of that pipeline currently can
and cannot make that trip.

\subsection{State of the field}\label{state-of-the-field}

\texttt{scikit-learn} (Pedregosa et al. 2011) defines the
\texttt{fit}/\texttt{transform} estimator contract that \texttt{Gators}
reuses directly, and its \texttt{preprocessing} module and
\texttt{ColumnTransformer} cover a subset of what \texttt{Gators}
implements, but operate on NumPy (Harris et al. 2020)/\texttt{pandas}
and are single-threaded by default. \texttt{feature-engine} (Galli 2021)
extends \texttt{scikit-learn} with a broader catalogue of encoders,
imputers, and discretizers on the same \texttt{pandas}-based footing.
\texttt{category\_encoders} (McGinnis et al. 2018) offers a further,
encoder-focused catalogue with the same estimator contract.
\texttt{Gators} overlaps in scope with all three but differs in its
execution engine --- \texttt{Polars}, a Rust-implemented (Matsakis and
Klock 2014), multi-threaded DataFrame library built on the Apache Arrow
columnar format (Apache Software Foundation 2016) with a lazy query
optimizer (Vink and Polars Contributors 2020) --- and in shipping a
built-in export path to a portable inference graph.
\texttt{sklearn-onnx} (ONNX Contributors 2019b) converts individual
\texttt{scikit-learn} estimators to ONNX one at a time; \texttt{Gators}
instead exports an entire fitted pipeline as one self-contained,
checker-validated ONNX \texttt{ModelProto}. TensorFlow Transform (Baylor
et al. 2017) takes the opposite approach to training/serving
consistency, compiling preprocessing directly into a TensorFlow serving
graph; \texttt{Gators} keeps training-time execution framework-agnostic
(\texttt{Polars}) and generates a portable graph only at export time, at
the cost of a documented, principled boundary on which transformers can
currently make that trip (Section ``Software design'').

\subsection{Software design}\label{software-design}

Every transformer subclasses a common \texttt{\_BaseTransformer},
combining Pydantic's (Colvin and Pydantic Contributors 2017)
\texttt{BaseModel} (typed, validated, keyword-only configuration) with
\texttt{scikit-learn}'s \texttt{BaseEstimator}/\texttt{TransformerMixin}
(\texttt{get\_params}, \texttt{set\_params}, pipeline compatibility).
\texttt{fit(X,\ y=None)} computes and stores statistics as private
attributes and returns \texttt{self}; \texttt{transform(X)} applies
those statistics and returns a new \texttt{Polars} DataFrame. A shared
\texttt{\_\_init\_subclass\_\_} hook gives every transformer consistent
\texttt{LazyFrame} materialization, a \texttt{NotFittedError} guard, and
automatic input-column/dtype bookkeeping, so that 107 independently
authored transformers share one calling convention without
reimplementing it individually. Internally, transformers batch
column-wise operations into a single Polars expression list per
\texttt{with\_columns}/\texttt{select} call rather than looping and
re-assigning the DataFrame column by column, which is the direct
mechanism behind the library's measured performance advantage over
\texttt{pandas}-based tooling: it lets one \texttt{transform} call
compile to one query plan that the engine can parallelize across columns
and cores, instead of many sequential ones.
\texttt{gators.pipeline.Pipeline} chains named steps exactly as
\texttt{sklearn.pipeline.Pipeline} does, remaining a valid
\texttt{scikit-learn} estimator to downstream tooling.
\texttt{gators.onnx\_converters.pipeline\_to\_onnx} walks a fitted
pipeline and, for each step, invokes a per-transformer-family ONNX
converter; the resulting graph is opset-selected automatically and
validated with \texttt{onnx.checker} before being returned. All
data-cleaning, clipping, encoding, imputation, discretization, scaling,
and numeric/datetime feature-generation transformers currently convert.
Nine string-feature transformers do not -- seven because they rely on
variable-length tokenization, fuzzy edit distance, or
regex-group/list-aggregate semantics that ONNX's string-tensor operator
set cannot express, and two because they require a string-length or
string-slice operator absent from the standard opset -- and the six
feature-selection transformers are not yet wired into the exporter.
\texttt{check\_pipeline\_onnx\_compatibility} lets a user audit a
pipeline for this boundary \emph{before} attempting a production export,
and the exporter itself supports both fail-fast and
pass-through-unchanged behavior for unsupported steps, rather than
failing opaquely.

\subsection{Research impact statement}\label{research-impact-statement}

\texttt{Gators} was developed by PayPal's Payment Service Provider (PSP)
Data Team to support fraud-detection and risk-modeling workflows, where
preprocessing and feature-engineering cost is a first-order concern at
the row counts typical of transaction data, and where a consistent
training-to-serving path materially reduces the risk of silent model
degradation from re-implementation drift. Its open, rerunnable benchmark
suite (comparing \texttt{Gators} against \texttt{scikit-learn} and
\texttt{feature-engine} across row counts, feature counts, missingness,
cardinality, and thread configurations, and comparing native execution
against the ONNX-exported path across batch sizes) gives downstream
researchers a reproducible reference point for preprocessing performance
claims in tabular ML pipelines, rather than a single, uncontextualized
number. By documenting precisely which transformer families can and
cannot currently be represented in a portable inference graph, and why,
the project also contributes a concrete characterization of where
graph-based, statically-shaped inference formats reach their limits for
tabular preprocessing --- a boundary condition relevant to any research
building ONNX-based deployment tooling for classical (non-neural) ML
pipelines, not only to this library's own users. \texttt{Gators} is
available at \url{https://github.com/paypal/Gators/}.

\subsection{AI usage disclosure}\label{ai-usage-disclosure}

This paper was initially drafted by the authors. Portions of this
project's development were carried out with the assistance of an AI
coding agent (GitHub Copilot, built on large language models), operating
under direct human review and direction. This included: authoring and
running the extended benchmark suite
(\texttt{benchmarks/run\_matrix\_benchmarks.py},
\texttt{benchmarks/run\_onnx\_benchmarks.py}); and helping expand and
revise sections of the authors' draft, including this paper. All
AI-assisted code changes were reviewed, tested against the existing
2,323-test suite, and benchmarked by a human author before being
retained.

\subsection{Acknowledgements}\label{acknowledgements}

We thank Prem Thangamani for his managerial support of the development
of \texttt{Gators}, PayPal's PSP Data Team for supporting the project
more broadly, and the maintainers of \texttt{Polars},
\texttt{scikit-learn}, \texttt{Pydantic}, and \texttt{ONNX}, whose
libraries this project builds directly on.

\protect\phantomsection\label{refs}
\begin{CSLReferences}{1}{1}
\bibitem[\citeproctext]{ref-arrow}
Apache Software Foundation. 2016. \emph{Apache Arrow: A Cross-Language
Development Platform for in-Memory Data}.
\href{https://arrow.apache.org}{Https://arrow.apache.org}.

\bibitem[\citeproctext]{ref-baylor2017}
Baylor, Denis, Eric Breck, Heng-Tze Cheng, et al. 2017. {``{TFX}: A
{T}ensor{F}low-Based Production-Scale Machine Learning Platform.''}
\emph{Proceedings of the 23rd ACM SIGKDD International Conference on
Knowledge Discovery and Data Mining (KDD)}, 1387--95.
\url{https://doi.org/10.1145/3097983.3098021}.

\bibitem[\citeproctext]{ref-breck2017}
Breck, Eric, Shanqing Cai, Eric Nielsen, Michael Salib, and D. Sculley.
2017. {``The {ML} Test Score: A Rubric for {ML} Production Readiness and
Technical Debt Reduction.''} \emph{IEEE International Conference on Big
Data (Big Data)}, 1123--32.
\url{https://doi.org/10.1109/BigData.2017.8258038}.

\bibitem[\citeproctext]{ref-pydantic}
Colvin, Samuel, and Pydantic Contributors. 2017. \emph{Pydantic: Data
Validation Using {P}ython Type Hints}.
\href{https://docs.pydantic.dev}{Https://docs.pydantic.dev}.

\bibitem[\citeproctext]{ref-galli2021}
Galli, Soledad. 2021. {``Feature\_engine: A {P}ython Package for Feature
Engineering for Machine Learning.''} \emph{Journal of Open Source
Software} 6 (65): 3642. \url{https://doi.org/10.21105/joss.03642}.

\bibitem[\citeproctext]{ref-harris2020numpy}
Harris, Charles R., K. Jarrod Millman, Stéfan J. van der Walt, et al.
2020. {``Array Programming with {NumPy}.''} \emph{Nature} 585: 357--62.
\url{https://doi.org/10.1038/s41586-020-2649-2}.

\bibitem[\citeproctext]{ref-matsakis2014rust}
Matsakis, Nicholas D., and Felix S. Klock. 2014. {``The {R}ust
Language.''} \emph{ACM SIGAda Ada Letters} 34 (3): 103--4.
\url{https://doi.org/10.1145/2692956.2663188}.

\bibitem[\citeproctext]{ref-category_encoders}
McGinnis, Will, Chapman Siu, Shea Andre, and Hanyu Huang. 2018.
\emph{Category\_encoders: A Scikit-Learn-Contrib Package of Transformers
for Encoding Categorical Data}.
\href{https://github.com/scikit-learn-contrib/category_encoders}{Https://github.com/scikit-learn-contrib/category\_encoders}.

\bibitem[\citeproctext]{ref-onnx}
ONNX Contributors. 2019a. \emph{{ONNX}: Open Neural Network Exchange}.
\href{https://onnx.ai}{Https://onnx.ai}.

\bibitem[\citeproctext]{ref-sklearnonnx}
ONNX Contributors. 2019b. \emph{Sklearn-Onnx: Convert Scikit-Learn
Models to {ONNX}}.
\href{https://github.com/onnx/sklearn-onnx}{Https://github.com/onnx/sklearn-onnx}.

\bibitem[\citeproctext]{ref-pedregosa2011}
Pedregosa, Fabian, Gaël Varoquaux, Alexandre Gramfort, et al. 2011.
{``Scikit-Learn: Machine Learning in {P}ython.''} \emph{Journal of
Machine Learning Research} 12: 2825--30.

\bibitem[\citeproctext]{ref-prokhorenkova2018}
Prokhorenkova, Liudmila, Gleb Gusev, Aleksandr Vorobev, Anna Veronika
Dorogush, and Andrey Gulin. 2018. {``{CatBoost}: Unbiased Boosting with
Categorical Features.''} \emph{Advances in Neural Information Processing
Systems (NeurIPS)} 31.

\bibitem[\citeproctext]{ref-sculley2015}
Sculley, D., Gary Holt, Daniel Golovin, et al. 2015. {``Hidden Technical
Debt in Machine Learning Systems.''} \emph{Advances in Neural
Information Processing Systems (NeurIPS)} 28.

\bibitem[\citeproctext]{ref-polars}
Vink, Ritchie, and Polars Contributors. 2020. \emph{Polars: Fast
Multi-Threaded, Hybrid-Out-of-Core {D}ata{F}rame Library}.
\href{https://www.pola.rs}{Https://www.pola.rs}.

\end{CSLReferences}

\end{document}
